\scanpage{28}
\begin{center}$-\ 32\ -$\end{center}
\hfill n$^\circ$ 371
\[
\rho_*(F)\otimes_A E\longrightarrow\rho_*(F\otimes_B\rho^*(E))
\]
(chap. II, $3^{\mathrm e}$ éd., \S\ 3, n$^\circ$ 2, cor. de la prop. 6) \uline{est un isomorphisme de} $A$-\uline{algèbres}.

Les vérifications sont triviales.

En particulier, la structure de $A$-algèbre \unsure{de} $E_{(B)}=B\otimes_A E$ obtenue par restriction à $A$ de l'anneau des scalaires, est identique à la structure de l'algèbre $B\otimes_A E$, produit tensoriel des $A$-algèbres $B$ et $E$.

\uline{Exemples de produits tensoriels d'algèbres.}--- 1) Soient $A$ un anneau commutatif, $M,N$ deux $A$-modules ; l'application canonique
\[
\operatorname{End}_A(M)\otimes_A\operatorname{End}_A(N)\longrightarrow\operatorname{End}_A(M\otimes_A N)
\]
(chap. II, $3^{\mathrm e}$ éd., \S\ 4, n$^\circ$ 4) est un homomorphisme de $A$-\uline{algèbres}, comme il résulte du chap. II, $3^{\mathrm e}$ éd., \S\ 3, n$^\circ$ 2, formule (5). Lorsque $M$ ou $N$ est un $A$-module \uline{projectif de type fini}, on sait que cet homomorphisme est \uline{bijectif} (chap. II, $3^{\mathrm e}$ éd., \S\ 4, n$^\circ$ 4, prop. 4). En particulier, on a un isomorphisme canonique d'\uline{algèbres de matrices} sur $A$
\[
M_r(A)\otimes_A M_s(A)\longrightarrow M_{rs}(A).
\]

2) Soient $S,T$ deux monoïdes quelconques, $E=A^{(S)}$ et $F=A^{(T)}$ les algèbres des monoïdes $S$ et $T$ relatives à l'anneau $A$ (\S\ 1, n$^\circ$ 11) ; on a alors un isomorphisme canonique du produit tensoriel $E\otimes_A F$ des algèbres $E$ et $F$ sur l'\uline{algèbre} $G=A^{(S\times T)}$ \uline{du monoïde produit} $S\times T$ relative à $A$. En effet, les éléments $u\otimes v$ (resp. $(u,v)$), où $u$ parcourt $S$ et $v$ parcourt $T$, forment une base de $E\otimes_A F$ en vertu du chap. II, $3^{\mathrm e}$ éd., \S\ 3, n$^\circ$ 7 (resp. de $G$) ; l'isomorphisme considéré s'obtient en faisant correspondre $(u,v)$ à $u\otimes v$ et il résulte immédiatement des définitions que c'est bien un isomorphisme d'\uline{algèbres}.

2. \uline{Caractérisation universelle des produits tensoriels d'algèbres.}

Nous reprenons nos conventions générales, les algèbres considérées étant
\mspage{27}
